\documentclass[conference]{IEEEtran}
\IEEEoverridecommandlockouts

% Packages
\usepackage{cite}
\usepackage{amsmath,amssymb,amsfonts}
\usepackage{algorithmic}
\usepackage{graphicx}
\usepackage{textcomp}
\usepackage{xcolor}
\usepackage{booktabs}
\usepackage{multirow}
\usepackage{tikz}
\usepackage{float}
\usetikzlibrary{shapes.geometric,arrows.meta,positioning}
\usepackage{pgfplots}
\pgfplotsset{compat=1.16}

% Corrected BibTeX definition
\def\BibTeX{{\rm B\kern-.05em{\sc i\kern-.025em b}\kern-.08em
    T\kern-.1667em\lower.7ex\hbox{E}\kern-.125emX}}

\begin{document}

\title{Typed-Decision Pre-Filtering for Cost-Efficient Academic Integrity Detection Pipelines}

\author{
\begin{tabular*}{\textwidth}{@{\extracolsep{\fill}}ccc}
\begin{minipage}[t]{0.30\textwidth}
\centering
\textbf{Shaik John Saida}\\
\textit{Department of Computer Science}\\
\textit{\& Engineering}\\
\textit{SRM University AP}\\
Neerukonda, Mangalagiri-522240\\
Guntur District, Andhra Pradesh, India\\
\small j9490490885@gmail.com
\end{minipage}
&
\begin{minipage}[t]{0.30\textwidth}
\centering
\textbf{Potti Sri Lakshmi Chetan}\\
\textit{Department of Computer Science}\\
\textit{\& Engineering}\\
\textit{SRM University AP}\\
Neerukonda, Mangalagiri-522240\\
Guntur District, Andhra Pradesh, India\\
\small pottisrilakshmichetanchetan@gmail.com
\end{minipage}
&
\begin{minipage}[t]{0.30\textwidth}
\centering
\textbf{Kothuri Pujith Sai Ram}\\
\textit{Department of Computer Science}\\
\textit{\& Engineering}\\
\textit{SRM University AP}\\
Neerukonda, Mangalagiri-522240\\
Guntur District, Andhra Pradesh, India\\
\small pujithsairamkothuri@gmail.com
\end{minipage}
\\
\end{tabular*}
}

\maketitle

% Abstract
\begin{abstract}
Large-scale academic integrity detection requires a balance between
reliable similarity assessment and computational efficiency. Lexical
methods are inexpensive but may miss paraphrased content, whereas
richer feature-based models require greater computation. This paper
presents a typed-decision pre-filtering approach that uses a lightweight
similarity model to identify straightforward cases and reserves a
comprehensive feature-based model for cases requiring further review.
The proposed feature representation extends an eight-feature baseline
with word-overlap, character n-gram, edit-distance, semantic, structural,
and phonetic features. The approach is evaluated on a text-plagiarism
benchmark (Dataset 1) and Quora Question Pairs (Dataset 2) using Matthews
Correlation Coefficient (MCC), accuracy, ROC-AUC, and inference latency. The complete
83-feature model improves MCC over the baseline from 0.3573 to 0.4705
on Dataset 1 and from 0.1990 to 0.3958 on Dataset 2. The lightweight baseline is faster
by factors of 2.8 and 20.7, respectively. These findings motivate a
two-stage architecture; however, end-to-end cascade performance remains
to be validated experimentally.
\end{abstract}

\begin{IEEEkeywords}
Academic integrity, plagiarism detection, paraphrase detection,
pre-filtering, feature engineering, Matthews correlation coefficient
\end{IEEEkeywords}

% =========================================================
% 1. INTRODUCTION
% =========================================================

\section{Introduction}

Checking for plagiarism is commonly used to ensure the upholding of
academic integrity. Methods such as transformer cross-encoders,
Abstract Syntax Trees or even large language models are computationally
intensive and so even more so when considering several institutions and
the comparisons that must be made between each submission~\cite{reimers2019}.
Although lexically based methods are typically low cost, they can fail
to recognise paraphrased material and/or semantic equivalence. Poor
similarity assessment can therefore result in additional, unwarranted

reviews. This has implications for both academic review panels and
students.
A solution is required that is both cost effective and provides a
reliable similarity assessment. An initial triage process can further

minimise cost by deciding which submission pairs require deeper
review. Here, we seek to create an improved, cost-effective
plagiarism detection solution. We wish to do this without solely
depending on low-cost lexical similarity.

In order to identify academic integrity concerns, we have created a
triage approach. Simple situations are addressed using a basic model
whereas more complicated instances are analysed by a more detailed
review model.

We enhance the eight features described in Sajid et al. (2025) \cite{sajid2025} by
incorporating lexical, character, edit distance, semantic,
structural and phonetic features. This extended set of features is
analysed on two data sets in order to assess both its classification
performance and cost effectiveness. Using our analysis of performance
and speed we suggest a pre-screening type system which categorises
plagiarism detection into two steps according to type.

% =========================================================
% 2. RELATED WORK
% =========================================================

\section{Related Work}

Plagiarism and paraphrase detection methods fall broadly into three
groups. \emph{Lexical methods} compare surface units such as words,
$n$-grams, and edit operations. TF-IDF cosine similarity with a
decision threshold is a common baseline \cite{amirzhanov2025}, and
hybrid overlap scores combine Jaccard similarity with normalised
overlap counts \cite{singh2026}.

\emph{Feature-based supervised methods}, to which the baseline of
Sajid et al. \cite{sajid2025} belongs, feed a compact set of
similarity features into a classical classifier. \emph{Neural
verifiers}, such as Sentence-BERT and cross-encoders \cite{reimers2019},
provide semantic discrimination but can have higher per-pair
computational costs.

Latent Semantic Analysis (LSA) \cite{deerwester1990} helps identify
similarities in meaning between texts, even when they use different
words. Levenshtein distance \cite{levenshtein1966} and Jaro--Winkler
similarity \cite{winkler1990} capture near-duplicates with small
edits, while Soundex \cite{russell1918} captures phonetic similarity.

Most prior work evaluates a complete feature set holistically;
the incremental value of each feature family and its dependence
on the dataset are rarely measured. MCC is used as the primary
metric because it is robust to class imbalance \cite{chicco2020}.

Dataset 1 is a text-plagiarism benchmark containing reference-suspicious
text pairs with plagiarism labels, using a validation/test split with
9,824 test pairs. Dataset 2 is Quora Question Pairs with paraphrase labels,
using 8,000 test pairs with a 5,000-pair training subset (from 49,401
available pairs) for computational tractability. These datasets are
described in detail in Section IV-A.

% =========================================================
% CONTRIBUTIONS
% =========================================================

\subsection{Contributions}

\begin{itemize}

    \item A similarity feature set including word overlap,
    character n-grams, edit distance, semantic, structural,
    and phonetic features, improving on the eight-feature
    baseline presented in Sajid et al. (2025) \cite{sajid2025}.

    \item Quantitative comparison of performance between the
    baseline and full feature set on two datasets. Performance
    is measured using classification metrics and inference speed.

    \item Qualitative comparison of the trade-offs between
    lightweight screening and thorough verification.

    \item A type-based pre-filtering system specifying how a
    lightweight model and a thorough model can be used in a
    two-stage detection process.

\end{itemize}

% =========================================================
% 3. METHODOLOGY
% =========================================================

\section{Methodology}

\subsection{Typed-Decision Pre-Filtering Architecture}

A flow is shown in Fig.~\ref{fig:proposed_architecture}.
All pairs are initially analysed using a basic filter (Models 1-2).
High certainty non-plagiaristic pairs are identified at this stage.
Lower certainty pairs are then passed onto a verification filter
(Model 3-5) and Model 5 represents the full system.
As with the review process, the majority of work is judged to be

non-plagiarised and so only a small number of items are submitted
for further analysis.
% ---------------------------------------------------------

% ARCHITECTURE IMAGE
% ---------------------------------------------------------

\begin{figure}[H]
\centering
\includegraphics[width=\linewidth]
{Gemini_Generated_Image_k17qifk17qifk17q.png}
\caption{Workflow of the proposed two-stage academic integrity
detection framework, integrating lightweight pre-filtering,
confidence-based routing, and comprehensive feature-based
verification.}
\label{fig:proposed_architecture}
\end{figure}
% ---------------------------------------------------------

% BASELINE FEATURES
% ---------------------------------------------------------

\subsection{Baseline Features (Model 1, 8 dimensions)}

The baseline of Sajid et al. \cite{sajid2025} includes
word Jaccard, character trigram Jaccard, number of overlapping
words, ratio of lengths, short-text overlap, exact match,

normalised overlap (based on larger text), and bigram

containment.
% ---------------------------------------------------------
% FEATURE CONTRIBUTIONS
% ---------------------------------------------------------
\subsection{Contributions C1--C4}
\noindent\textbf{C1: Word-Overlap Count Features (8 dimensions).}
This combines the word Jaccard, number of overlapping words, length
ratio, short-text overlap and word bigrams Jaccard along with
exact match. Normalised overlap and bigram containment provide
information on multi-word level similarities in addition to
single word level information.
\par\noindent\textbf{C2: Character N-Gram Features (5 dimensions).}
Character 2, 3, 4 and 5-grams Jaccard similarities, and trigram
containment, represent orthographic and morphological differences.
\par\noindent\textbf{C3: Edit Distance Features (4 dimensions).}
Word and character level Levenshtein similarities $1-d_{\mathrm{lev}}/\max(m,n)$ and
word and character level Jaro-Winkler similarities, for analysing
close matches that have experienced minor additions, omissions
or changes.
\par\noindent\textbf{C4: Semantic, Structural and Phonetic Features (58 dimensions).}
Latent topic semantics via LSA (25 dims per sentence; 50 dims total; formed using
a truncated singular value decomposition on a TF-IDF matrix with up to 5000 features),
question structure features (5 dims; common initial word, common terminal
word, both texts initiate with a question word, both texts
initiate with an auxiliary verb and common question words)
alongside Soundex phonetic features (3 dims; Jaccard similarity of
Soundex codes, number of common codes and code containment).

% =========================================================
% 4. EXPERIMENTAL SETUP
% =========================================================

\section{Experimental Setup}

\subsection{Datasets}

Two datasets are used under identical train/test splits:

\textbf{Dataset 1 --- Text Plagiarism Benchmark:} A repository text plagiarism
benchmark containing reference-suspicious text pairs with plagiarism labels.
The dataset uses a validation/test split with 9,824 test pairs. The \texttt{val}
split is used for training and the \texttt{test} split for evaluation.

\textbf{Dataset 2 --- Quora Question Pairs:} Question pairs from the Quora
dataset with paraphrase labels. For computational tractability, we use a
5,000-pair training subset (from 49,401 available pairs) and the full
8,000-pair test set.

% ---------------------------------------------------------
% PREPROCESSING
% ---------------------------------------------------------

\subsection{Preprocessing}

All models share identical preprocessing:
\begin{enumerate}
    \item Lowercase the text
    \item Remove non-alphanumeric characters (except spaces)
    \item Collapse multiple whitespace into single space
    \item Trim leading and trailing whitespace
\end{enumerate}

All models implement logistic regression with
\texttt{class\_weight=balanced} and \texttt{max\_iter=2000}
(random state 42).

The threshold used to classify pairs is selected to maximise
the MCC of the training data using a coarse search in the range
[0.05, 0.95], followed by a fine search in the vicinity of the
optimal value. The resulting threshold is applied directly to
the test data.

The metrics used are accuracy, precision, recall, F1, specificity,
ROC-AUC, PR-AUC, and MCC (primary metric). Model latency
(in seconds) is calculated as the average time taken to make
predictions for pairs in the test set.

The models are implemented using scikit-learn \cite{pedregosa2011}.

% =========================================================
% 5. RESULTS AND DISCUSSION
% =========================================================

\section{Results and Discussion}

\subsection{Quantitative Analysis}

The baseline model (Model 1) and the complete feature model
(Model 5) were evaluated on both datasets. Model 1 uses the
eight baseline features, whereas Model 5 uses the full
83-dimensional representation, including the additional
feature groups.

Table~\ref{tab:quantitative} summarises the principal
classification and latency results.

\begin{table}[!t]
\caption{Quantitative Comparison of the Baseline and Complete Feature Model}
\label{tab:quantitative}
\centering
\small
\setlength{\tabcolsep}{4pt}
\begin{tabular}{llcc}
\toprule
Dataset & Metric & Baseline (M1) & Proposed (M5) \\
\midrule
\multirow{4}{*}{Dataset 1}
& Accuracy & 0.7149 & 0.7641 \\
& MCC & 0.3573 & 0.4705 \\
& ROC-AUC & 0.7375 & 0.8007 \\
& Latency (ms/pair) & 0.062 & 0.176 \\
\midrule
\multirow{4}{*}{Dataset 2}
& Accuracy & 0.6009 & 0.7046 \\
& MCC & 0.1990 & 0.3958 \\
& ROC-AUC & 0.6367 & 0.7585 \\
& Latency (ms/pair) & 0.054 & 1.116 \\
\bottomrule
\end{tabular}
\end{table}

On Dataset 1, the complete model increases accuracy from 0.7149
to 0.7641, MCC from 0.3573 to 0.4705, and ROC-AUC from
0.7375 to 0.8007. The MCC improvement is approximately
31.7\% relative to the baseline.

On Dataset 2, accuracy increases from 0.6009 to 0.7046, MCC from
0.1990 to 0.3958, and ROC-AUC from 0.6367 to 0.7585,
corresponding to an approximately 98.9\% relative improvement
in MCC. Thus, the richer representation improves the reported
classification metrics on both datasets.

% ---------------------------------------------------------
% RESULT FIGURE 1
% ---------------------------------------------------------

\begin{figure}[!t]
    \centering
    \includegraphics[width=\linewidth]
    {Screenshot 2026-11-21 115958.png}
    \caption{Classification performance comparison between
    the baseline and proposed model on the evaluated datasets.}
    \label{fig:performance1}
\end{figure}

% ---------------------------------------------------------
% RESULT FIGURE 2
% ---------------------------------------------------------

\begin{figure}[!t]
    \centering
    \includegraphics[width=\linewidth]
    {Screenshot 2026-11-21 120043.png}
    \caption{Comparison of the classification performance
    metrics for the evaluated models.}
    \label{fig:performance2}
\end{figure}

The improvement is accompanied by increased inference latency.
On Dataset 1, the baseline requires 0.062 ms per pair, compared with
0.176 ms for the complete model; the baseline is approximately
2.8 times faster.

On Dataset 2, the corresponding latencies are 0.054 ms and 1.116 ms,
making the baseline approximately 20.7 times faster.

These results quantify the trade-off between classification
performance and computational cost. The measurements are
standalone model results and should not be interpreted as
measured end-to-end cascade performance.

% ---------------------------------------------------------
% QUALITATIVE ANALYSIS
% ---------------------------------------------------------

\subsection{Qualitative Analysis}

The baseline model offers a low-cost first stage because it
relies on a compact set of lexical and similarity features,
including word Jaccard, character-trigram Jaccard, overlap
count, length difference, exact-match indication, TF-IDF
cosine similarity, and LSA cosine similarity.

Its low latency makes it suitable for preliminary screening.
However, a compact feature representation may provide
insufficient evidence when two submissions express similar
meaning through substantial lexical changes.

The complete model extends the baseline with the four feature
families. Word-overlap features capture shared terms and
phrases; character $n$-grams capture local spelling and
morphological variation; edit-distance features identify
near-duplicate strings; and semantic, structural, and phonetic
features provide complementary evidence beyond exact lexical
overlap.

The higher MCC and ROC-AUC values observed for Model 5
indicate that this richer representation is more discriminative
on the two evaluated datasets. The size of the improvement
differs between Dataset 1 and Dataset 2, indicating that the benefit of
additional similarity evidence depends on dataset
characteristics.

The proposed architecture uses these complementary properties
in a staged workflow. The lightweight model is assigned to the
initial screening stage, while the comprehensive model is
intended to assess pairs that cannot be confidently resolved
in the first stage.

This arrangement is designed to avoid applying the more
expensive feature computation to every pair. A practical
deployment must define and validate the confidence thresholds
used for routing, and must ensure that potentially plagiarised
pairs are not incorrectly discarded by the first stage.

The current experiments establish the standalone performance
and latency of the two component models. They do not report
the fraction of pairs that would be resolved by the pre-filter,
the recall of the routed candidate set, or measured end-to-end
cascade latency.

Consequently, the architecture is a motivated deployment
design rather than a fully validated cascade result. These
quantities must be measured before claiming realised workload
reduction or end-to-end efficiency.

% ---------------------------------------------------------
% COST ANALYSIS
% ---------------------------------------------------------
\subsection{Cost Analysis}

The computational cost of the proposed model configurations
was evaluated by comparing execution latency and MCC per
millisecond for the lightweight Model 1 and the complete
Model 5.

On Dataset 1, Model 1 required 0.062 ms per pair and achieved an
MCC of 0.3573, resulting in an MCC-per-millisecond value
of 5.763. In comparison, Model 5 required 0.176 ms per pair,
achieved an MCC of 0.4705, and produced an
MCC-per-millisecond value of 2.673.

Thus, Model 1 was approximately 2.8 times faster than
Model 5 on Dataset 1.

% ---------------------------------------------------------
% COST FIGURE 1
% ---------------------------------------------------------

\begin{figure}[!t]
    \centering
    \includegraphics[width=\linewidth]
    {Screenshot 2026-11-21 120543.png}
    \caption{Comparison of model latency and MCC on Dataset 1.}
    \label{fig:cost1}
\end{figure}

% ---------------------------------------------------------
% COST FIGURE 2
% ---------------------------------------------------------

\begin{figure}[!t]
    \centering
    \includegraphics[width=\linewidth]
    {Screenshot 2026-11-21 120916.png}
    \caption{Comparison of model latency and MCC on Dataset 2.}
    \label{fig:cost2}
\end{figure}

On Dataset 2, Model 1 required 0.054 ms per pair and achieved
an MCC of 0.1990, resulting in an MCC-per-millisecond
value of 3.685.

Model 5 required 1.116 ms per pair and achieved an MCC
of 0.3958, producing an MCC-per-millisecond value of 0.355.
Consequently, Model 1 was approximately 20.7 times faster
than Model 5 on Dataset 2.

These measurements demonstrate the trade-off between
classification performance and computational cost. Although
Model 5 achieves higher MCC values on both datasets, its
execution time is also higher due to the additional feature
computations, including the LSA projection and edit-distance
operations.

This difference is particularly pronounced on Dataset 2, where
Model 5 requires substantially more execution time than
the lightweight baseline.

The measured latency values motivate the proposed two-stage
architecture, in which Model 1 acts as a
lightweight pre-filter and Model 5 performs comprehensive
verification on unresolved pairs.

Let $p$ represent the proportion of pairs resolved by
the pre-filter, $t_1$ represent the execution time of Model 1,
and $t_5$ represent the execution time of Model 5.

Assuming every pair is first processed by Model 1 and
only the remaining proportion $(1-p)$ is forwarded to
Model 5, the estimated average execution cost per pair is:

\begin{equation}
T_{\mathrm{cascade}} = t_1 + (1-p)t_5.
\label{eq:cascade_cost}
\end{equation}

For the cascade to have a lower average execution cost
than applying Model 5 to every pair, the following
condition must hold:

\begin{equation}
t_1 + (1-p)t_5 < t_5.
\end{equation}

Rearranging this expression gives the minimum pre-filter
resolution proportion:

\begin{equation}
p > \frac{t_1}{t_5}.
\end{equation}

Using the measured execution times, the estimated crossover
point is approximately 35.2\% for Dataset 1, based on the 0.062 ms
for Model 1 and 0.176 ms for Model 5.

For Dataset 2, the corresponding crossover point is approximately
4.8\%, based on 0.054 ms for Model 1 and 1.116 ms for
Model 5.

Therefore, the pre-filter would need to resolve more than
these respective proportions of pairs for the proposed
cascade to achieve a lower average execution cost than
using Model 5 alone.

However, these crossover values are theoretical estimates
derived from standalone model execution times. The actual
proportion of pairs resolved by the pre-filter was not
measured.

The end-to-end cascade latency, routing behaviour, and
classification performance therefore remain to be
experimentally evaluated.

% =========================================================
% 6. DEPLOYMENT RECOMMENDATION
% =========================================================

\section{Deployment Recommendation}

Based on the experimental results, Model 1 is a suitable
candidate for the lightweight pre-filtering stage because
of its lower execution latency on both datasets.

Under the proposed architecture, pairs confidently identified
as non-plagiarised could be resolved during the first stage,
while ambiguous pairs could be forwarded to Model 5 for
comprehensive verification.

The results support the potential use of the lightweight
pre-filter followed by the more comprehensive verification
stage. Nevertheless, this recommendation is based on the
standalone performance and latency measurements of the
individual models.

The complete cascade must be evaluated under different
routing thresholds to establish its actual efficiency and
classification performance.

% =========================================================
% 7. CONCLUSION
% =========================================================

\section{Conclusion}

This paper presented a typed-decision pre-filtering design
for cost-conscious academic integrity detection.

The complete 83-feature model improved MCC over the
eight-feature baseline by approximately 31.7\% on the
text-plagiarism benchmark and 98.9\% on Quora Question
Pairs.

It also achieved higher accuracy and ROC-AUC on both
datasets, while requiring greater per-pair inference time.

The measured latency difference motivates a two-stage
design in which a lightweight model performs initial
screening and a comprehensive model examines unresolved
cases.

However, the experiments reported here evaluate the
component models independently; end-to-end cascade
performance, candidate recall, routing rate, and realised
workload reduction remain to be measured.

Future work will implement and evaluate the complete
cascade, calibrate routing thresholds, assess candidate
recall and false-negative risk, and investigate
transformer-based verification and additional plagiarism
domains.

% =========================================================
% REFERENCES
% =========================================================

\begin{thebibliography}{00}

\bibitem{sajid2025}
Sajid, M., et al. ``Comparative analysis of text-based plagiarism
detection techniques.'' \emph{PLoS ONE}, vol. 20, no. 4,
e0319551, 2025. DOI: 10.1371/journal.pone.0319551.

\bibitem{amirzhanov2025}
Amirzhanov, A., et al. ``TF-IDF Cosine Similarity for Paraphrase
Detection.'' \emph{Frontiers in Computer Science}, 2025.

\bibitem{singh2026}
Singh, S. and Gupta, R. ``Hybrid Word-Overlap Approach for
Paraphrase Detection.'' \emph{Knowledge and Information Systems},
2026.

\bibitem{chicco2020}
D. Chicco and G. Jurman, ``The advantages of the Matthews
correlation coefficient (MCC) over F1 score and accuracy in
binary classification evaluation.'' \emph{BMC Genomics},
vol. 21, art. 6, 2020.

\bibitem{deerwester1990}
S. Deerwester, S. T. Dumais, G. W. Furnas, T. K. Landauer,
and R. Harshman, ``Indexing by latent semantic analysis,''
\emph{J. Amer. Soc. Inf. Sci.}, vol. 41, no. 6,
pp. 391--407, 1990.

\bibitem{levenshtein1966}
V. I. Levenshtein, ``Binary codes capable of correcting
deletions, insertions, and reversals,'' \emph{Soviet Physics
Doklady}, vol. 10, no. 8, pp. 707--710, 1966.

\bibitem{winkler1990}
W. E. Winkler, ``String comparator metrics and enhanced
decision rules in the Fellegi--Sunter model of record
linkage,'' in \emph{Proc. Section on Survey Research
Methods, ASA}, 1990, pp. 354--359.

\bibitem{russell1918}
R. C. Russell, ``Index,'' U.S. Patent 1,261,167, 1918.

\bibitem{reimers2019}
N. Reimers and I. Gurevych, ``Sentence-BERT: Sentence
embeddings using Siamese BERT-networks,'' in
\emph{Proc. EMNLP-IJCNLP}, 2019, pp. 3982--3992.

\bibitem{pedregosa2011}
F. Pedregosa et al., ``Scikit-learn: Machine learning
in Python,'' \emph{J. Mach. Learn. Res.}, vol. 12,
pp. 2825--2830, 2011.

\end{thebibliography}

\end{document}